% !PRINTED 458
\weberpage{458}
The theorem is first of all evident when the functions $\Phi_1, \Phi_2, \Phi_3\ldots$ depend on only one variable; for then there are altogether only a finite number of numerical values for this variable for which one of these functions becomes zero.

But then we can easily see the correctness of the theorem for functions of $n+1$ variables, provided we regard it as proved for functions of $n$ variables. For if we order the functions according to the $(n+1)$th variable $t$, then by hypothesis we can put for the remaining $n$ variables such rational values that in none of the functions do all the coefficients of the powers of $t$ vanish; we then have functions of the single variable $t$, and for it we can put such a rational value that none of the functions vanishes. The theorem is thus proved in general.

It shall now first be shown that the simultaneous adjunction of \gesperrt{several} algebraic quantities can be replaced by the adjunction of \gesperrt{a single one}; in other words, that every field $\Omega(\alpha, \beta, \gamma\ldots)$ can be regarded as a field $\Omega(\alpha)$.

Let then $\alpha, \beta, \gamma\ldots$ be algebraic quantities in arbitrary number, and let

\weq{1}{A(x) = 0, \quad B(x) = 0, \quad C(x) = 0\ldots}

be equations in $\Omega$ whose roots are $x=\alpha$, $x=\beta$, $x=\gamma\ldots$. None of these equations is to have a multiple root, an assumption by which the generality is not restricted.

We have seen in the preceding section that every quantity of a field $\Omega(\alpha)$ can be represented as an \gesperrt{integral rational function} in $\Omega$ of $\alpha$. If in this theorem we put in place of the field $\Omega$ the field $\Omega(\beta, \gamma\ldots)$, it follows that every number in $\Omega(\alpha, \beta, \gamma\ldots)$ can be represented as an integral rational function of $\alpha$ with coefficients from $\Omega(\beta, \gamma\ldots)$; and if we repeat the same mode of inference with respect to the coefficients, we obtain:

2. \gesperrt{Every quantity of the field $\Omega(\alpha, \beta, \gamma\ldots)$ can be represented as an integral rational function in $\Omega$ of $\alpha, \beta, \gamma\ldots$.}
